%%%%%%%%%%%%%%%%%%%%%%%%%%%%%%%%%%%%%%%%%%%%%%%%%%%%%%%%%%%
% PREAMBOLO
%%%%%%%%%%%%%%%%%%%%%%%%%%%%%%%%%%%%%%%%%%%%%%%%%%%%%%%%%%%

 % === Impostazione del documento ==========================
\documentclass[12pt,a4paper,oneside,italian,hidelinks]{book}
\usepackage{setspace}
\onehalfspace

% === Regolazione dei margini =============================\\
%\addtolength{\oddsidemargin}{30pt}
%\addtolength{\evensidemargin}{-30pt}
\usepackage{fancyhdr}
\usepackage{multirow}
\usepackage{rotating}
\usepackage{multicol}
\usepackage[section]{placeins}
\usepackage[usenames,dvipsnames]{color}
\usepackage{appendix}
\usepackage{chngpage}
\usepackage[margin=3cm]{geometry}

% === Impostazione dei font ===============================
\usepackage[utf8x]{inputenc}
\usepackage[italian]{babel}
\usepackage[10pt]{type1ec}
\usepackage[T1]{fontenc}
\usepackage{ae}
\usepackage{relsize}
\usepackage{csquotes} 
\usepackage{amsmath}
\usepackage{amsfonts}
\usepackage{mathdots}
\usepackage[colorlinks=true]{hyperref}
\hypersetup{
	bookmarksnumbered=true,
	linkcolor=myblue,
	citecolor=black,
	urlcolor=black,
}
\usepackage{verbatim}
\usepackage{alltt}

% === Integrazione delle figure ===========================
\usepackage{graphicx}
\graphicspath{{./imgs/}}
\renewcommand{\figurename}{Fig.}
\usepackage{subfigure}
\usepackage{epstopdf}
\usepackage{caption}

% === Per gli algoritmi ===================================
\usepackage{algorithmicx}
\usepackage[ruled]{algorithm}
\usepackage{algpseudocode}
\usepackage{listings}

% === Per le tabelle ======================================
\usepackage{tabularx}
\usepackage{booktabs}

% === Per gli snippet =====================================
\definecolor{mygreen}{rgb}{0,0.6,0}
\definecolor{mygray}{rgb}{0.5,0.5,0.5}
\definecolor{mymauve}{rgb}{0.58,0,0.82}
\definecolor{myblue}{rgb}{0.030,0.237,0.443}
\lstset{
  backgroundcolor=\color{white},   % choose the background color; you must add \usepackage{color} or \usepackage{xcolor}
  basicstyle=\footnotesize,        % the size of the fonts that are used for the code
  breakatwhitespace=false,         % sets if automatic breaks should only happen at whitespace
  breaklines=true,                 % sets automatic line breaking
  captionpos=n,                    % sets the caption-position to bottom
  commentstyle=\color{mygreen},    % comment style
  deletekeywords={},			   % if you want to delete keywords from the given language
  escapeinside={\%*}{*)},          % if you want to add LaTeX within your code
  extendedchars=true,              % lets you use non-ASCII characters; for 8-bits encodings only, does not work with UTF-8
  frame=single,                    % adds a frame around the code
  keepspaces=true,                 % keeps spaces in text, useful for keeping indentation of code (possibly needs columns=flexible)
  keywordstyle=\color{blue},       % keyword style
  language=Java,                   % the language of the code
  morekeywords={},            	   % if you want to add more keywords to the set
  numbers=left,                    % where to put the line-numbers; possible values are (none, left, right)
  numbersep=5pt,                   % how far the line-numbers are from the code
  numberstyle=\tiny\color{mygray}, % the style that is used for the line-numbers
  rulecolor=\color{black},         % if not set, the frame-color may be changed on line-breaks within not-black text (e.g. comments (green here))
  showspaces=false,                % show spaces everywhere adding particular underscores; it overrides 'showstringspaces'
  showstringspaces=false,          % underline spaces within strings only
  showtabs=false,                  % show tabs within strings adding particular underscores
  stepnumber=1,                    % the step between two line-numbers. If it's 1, each line will be numbered
  stringstyle=\color{mymauve},     % string literal style
  tabsize=2,                       % sets default tabsize to 2 spaces
  title=\lstname                   % show the filename of files included with \lstinputlisting; also try caption instead of title
}

% === Per altre opzioni =====================================
\usepackage[italian]{varioref}
\usepackage[dvipsnames,svgnames,x11names]{xcolor}
\usepackage[hyphenbreaks]{breakurl}
\usepackage{lipsum}
\usepackage{calc}
\usepackage{enumitem}

\newlength{\useCaseLeft}
\newlength{\useCaseRight}
\newlength{\useCaseNum}
\newlength{\useCaseMulticol}
\newlength{\useCaseOneCol}
\newlength{\useCaseTwoCol}
\newlength{\useCaseThreeCol}
\newlength{\useCaseFourCol}
\newlength{\useCaseFiveCol}
\setlength{\useCaseLeft}{0.20\textwidth}
\setlength{\useCaseRight}{0.12\textwidth}
\setlength{\useCaseNum}{0.05\textwidth}
\setlength{\useCaseMulticol}{0.65\textwidth + 10\tabcolsep}
\setlength{\useCaseOneCol}{0.60\textwidth + 8\tabcolsep}
\setlength{\useCaseTwoCol}{0.30\textwidth + 3\tabcolsep}
\setlength{\useCaseThreeCol}{0.20\textwidth + 1.333\tabcolsep}
\setlength{\useCaseFourCol}{0.15\textwidth + 0.5\tabcolsep}
\setlength{\useCaseFiveCol}{0.12\textwidth}

%%%%%%%%%%%%%%%%%%%%%%%%%%%%%%%%%%%%%%%%%%%%%%%%%%%%%%%%%%%
% TESTO DEL RAD
%%%%%%%%%%%%%%%%%%%%%%%%%%%%%%%%%%%%%%%%%%%%%%%%%%%%%%%%%%%
\begin{document}
	% === Frontespizio ====================================
	\frontmatter
	\pagestyle{empty}
	\label{key}%%%%%%%%%%%%%%%%%%%%%%%%%%%%%%%%%%%%%%%%%%%%%%%%%%%%%%%%%%%
% Frontespizio

% vspace serve ad aggiungere extra spazio verticale
% em sta ad indicare la grandezza della lettera M maiuscola

% Large indica una dimensione del font di 14.4 pt
% large indica una dimensione del font di 12 pt
% normalsize indica una dimensione del font di 10 pt

% vfill inserisce sufficiente spazio binaco verticalmente per fare in modo che il
% sopra e il sotto del testo siano allieneati col margine superiore e inferiore
\begin{titlepage}
 \begin{center}
	 \includegraphics[width=3.5cm]{unimol/Unimol.eps}\\
	 \vspace{0.8em}
	 {\Large \textsc{Università degli studi del Molise}}\\
	 \vspace{0.8em}
	 {\Large \textsc{Dipartimento di Bioscienze e Territorio}}\\
	 \vspace{2em}
	 {\normalsize Corsi di Studio in}\\
	 \vspace{1em}
	 {\Large \textsc{Informatica e Sicurezza dei sistemi software}} \\
	 \vspace{1em}
	 {\normalsize Corsi di \textit{Ingegneria del Software} e  \textit{Gestione progetti software}}\\
	 \vspace{5em}
	{\LARGE\textbf{Applicazione \textit{Studenti Unimol}}:} \\
	\vspace{1em}
	{\LARGE\textbf{Requirement Analysis Document (RAD)}} \\
	\vspace{6em}
 \end{center}

\vskip 0.2cm
  \begin{center}
	\begin{tabular}{l c c c c c c l}
	  Docente & & & & & & & Versione 3.0\\[0.2cm]
	  \large{Chiar.mo Prof. Fausto Fasano} & & & & & & & 31/05/2019\\[1cm]
	\end{tabular}
%  \end{center}
\vfill
%\begin{center}
{\normalsize Anno Accademico 2018/2019}
\end{center}
\end{titlepage}
\clearpage
	\newpage
	\thispagestyle{empty}
	\mbox{}

	% === Indice ==========================================
	\setcounter{secnumdepth}{3}
	\setcounter{tocdepth}{3}
	\tableofcontents

	% === Capitoli RAD ===================================
	\pagestyle{plain}
	\newpage
	\mainmatter
	%%%%%%%%%%%%%%%%%%%%%%%%%%%%%%%%%%%%%h%%%%%%%%%%%%%%%%%%%%%%

\chapter{Dominio del problema}
\label{ref:Introduzione}
Sezione comune a tutti i gruppi: la curerà qualcuno di noi. Non rivolta ai singoli gruppi.
%%%%%%%%%%%%%%%%%%%%%%%%%%%%%%%%%%%%%%%%%%%%%%%%%%%%%%%%%%%

\section{Introduzione al RAD}

\paragraph{}
Questo \textit{Requirement Analysis Document} ha lo scopo di descrivere la fase di raccolta e analisi dei requisiti funzionali, non funzionali e pseudo-requisiti dell'applicazione \textit{Studenti Unimol} quale sistema da sviluppare per gli esami di \textit{Ingegneria del software e laboratorio} e \textit{Gestione progetti software} previsti dai piani di studio del \textit{Corso di studio unificato in Informatica}. Tale documento funge da contratto tra il prof. Fausto Fasano che ha commissionato il sistema e i team di progettazione e sviluppo che lo realizzeranno. Esso fornisce una panoramica astratta del sistema che gli studenti si accingeranno ad implementare.

\section{Scope}

\paragraph{}
Questo progetto ha come scopo quello di realizzare una App ufficiale per l’Università degli Studi del Molise. L’obiettivo è quello di rendere disponibile, agli studenti universitari, i principali servizi di Ateneo attraverso un’applicazione completa ed esaustiva.
In primo luogo, dal punto di vista informativo ed organizzativo, il sistema deve essere capace di focalizzarsi sull’esigenza comune agli studenti di accedere velocemente alle informazioni e ai dati afferenti al proprio Piano di Studio. In tal modo, lo studente può monitorare in modo agevole ed intuitivo tutto ciò che attiene al proprio percorso universitario e alle principali attività svolte all'interno dell'Ateneo.

La tecnologia usata facilita estremamente l'utilizzo dell'applicazione. Si tratta di un’app ibrida multipiattaforma, che permette di avere alte performance su ciascun dispositivo mobile, in particolar modo su smartphone \textit{iOS} e \textit{Android}.  L’app supporta funzionalità di base ma anche funzionalità avanzate, le quali garantiscono un servizio aggiuntivo per gli studenti che decidono di utilizzarla. Lo studente può visualizzare tutti i suoi dati universitari e compiere una serie di attività, accedendo con le sue credenziali del portale \textit{Esse3} alla propria area personale. Tra le funzionalità avanzate, introdotta nell'anno accademico 2018/2019, è presente un servizio di messagistica instantaneo.
L'idea di base è quella di fornire un sistema di comunicazione semplice ed affidabile, evitando la trasmissione di informazioni personali, quali i contatti personali, come numero di telefono ed email, di docenti e studenti.
Le chat sono create automaticamente in base alla \textit{Coorte accademica} e ai relativi \textit{corsi} e sono manutenute per un periodo limitato di tempo. La chat verrà abilitata per ciascuno studente o docente nel momento in cui egli fa il login nell’\textit{app}.


\section{Contesto e panoramica del sistema}
Il corso di \textit{Ingegneria del software e laboratorio} all'\textit{Università degli Studi del Molise} prevede la suddivisione degli studenti in gruppi di lavoro ai quali è chiesto di progettare, documentare e sviluppare alcune funzionalità interne all'applicazione \textit{Studenti Unimol} utilizzata dagli studenti per la gestione semplificata della loro carriera accademica. Sarà rilasciata una versione aggiornata di quella attuale previa revisione della versione già esistente. Essa sarà resa disponibile a tutti gli studenti regolarmente iscritti all'\textit{Università degli Studi del Molise}.

\section{Manager di progetto e sviluppatori}

\paragraph{}
L'app \textit{Studenti Unimol} è stata documentata e implementata nel 2017 dagli studenti del corso di \textit{Ingegneria del Software}, coordinati da gruppi di manager del corso magistrale in \textit{Sicurezza dei sistemi software}, operanti nell'ambito dell'insegnamento di \text{Gestione progetti software}. Successivamente l'applicazione è stata manutenuta dal prof. Fausto Fasano, il quale, nel 2019, ha chiesto agli studenti di triennale e magistrale di manutenere e far evolvere l'attuale applicazione per produrre la versione 3.0. Ogni team è supervisionato da alcuni manager che si occuperanno della comunicazione e dell'organizzazione dell'intero progetto, coordinando i gruppi a loro assegnati e gestendo le relazioni orizzontali tra i gruppi che lavorano in parallelo su funzionalità diverse dell'app. \\

\textbf {Gruppo 1 - Funzionalità: piano di studi, appelli, materiale didattico} \\ \\
\textbf{Manager} \\
Fantini Martina (PM), Fierro Fabiana (PM), Varriano Giulia (QM), Mastropaolo Antonio (SM). \\
\textbf{Ingegneri del software} \\
Caserio Walter, Ciaramella Giovanni, Daniele Raffaele, De Turris Antonio, Discenza Christian, Fagnano Stefano, Iannotti Carmine,  Muccigrosso Marco, Russodivito Marco, Tata Giancarlo. \\

\textbf{Gruppo 2 - Funzionalità: gestione notifiche e gestione orario} \\ \\
\textbf{Manager} \\
Piedimonte Massimo (PM), Placella Davide (PM), La Rocca Piera Elena (QM), Polisena Alessandro Bruno (QM), Di Tommaso Fabio (SM). \\
\textbf{Ingegneri del software} \\
Armenti Carmen, Buro Martina, Cocozza Nicola, Discenza Silvia, Lucchetti Roberto, Mazzocco Giuseppina, Schiavone Raffaele, Siravo Luca, Spina Chiara, Varrati Angelo Gino. \\

\textbf {Gruppo 3 - Funzionalità: previsione della media} \\ \\
\textbf{Manager} \\Crincoli Giuseppe (PM), Fabrizio Emilio (PM), D'Ercole Giovanna (QM), Di Cristino Gianluca (SM). \\
\textbf{Ingegneri del software} \\Cancelliere Alessandro, Di Pilla Francesco, Maglioli Raffaele, Marzullo Alessio, Placella Andrea, Stefanelli Vincenzo, Venditti Giorgio.\\

\textbf {Gruppo 4 - Funzionalità: news} \\ \\
\textbf{Manager} \\
Fausto Fasano. \\ \\ \\ \\

\textbf {Gruppo 5 - Funzionalità: rubrica} \\ \\
\textbf{Manager} \\
Carnevale Filippo (PM), Marinaro Tiziano (QM), Tortola Domenico (SM). \\
\textbf{Ingegneri del software} \\
Ciardiello Gabriele, Pizzi Mario, Vitiello Giò. \\

\textbf{Gruppo 6 - Funzionalità: chat per studenti e docenti con annesso pannello di amministrazione} \\ \\
\textbf{Manager} \\
\textbf{} 
Giovanni Rosa (PM), Michele Guerra (QM), Angelo Iallonardi (SM). \\
\textbf{Ingegneri del software } \\
Daniele Albanese, Antonio Antenucci, Mattia Ciccaglione, Andrea D'Aguanno, Antonio De Santis,  Francesco Di Rito, Gianluca Farinaro, Antonio Fratianni, Emanuela Guglielmi, Gaia Iannone, Aldo Palombo, Marica Principe, Federico Zappone.

\clearpage
\section{Glossario}

\paragraph{}
Di seguito è riportata una tabella degli acronimi e dei termini tecnici utilizzati nel documento:

\begin{table}[!h]
\begin{tabular}{p{1.5in}|p{4in}} \\
	{\bf Termine o sigla} & {\bf Descrizione} \\ \hline
	\textbf{RAD} & \textit{Requirement Analysis Document}, documento di analisi dei requisiti. \\
	\textbf{Attore} &  Entità esterna al sistema che interagisce con
	esso. In questo documento la definizione degli attori viene trattata al paragrafo \ref{sec:attori} \\
	\textbf{Modello di workflow} & Modello del flusso di lavoro: indica il modello di sistema che viene utilizzato dagli sviluppatori durante la fase di analisi dei requisiti e implementazione del sistema. \\
	\textbf{Scenario} & Singola istanza di un caso d’uso. Descrive concretamente una situazione ipotetica che potrebbe verificarsi all'avverarsi di un caso d’uso. \\
	\textbf{Caso d'uso} & Flusso di eventi che coinvolge il sistema e alcuni attori che si genera quando  un attore interagisce con il sistema e, poiché sono soddisfatte delle condizioni di ingresso e dei vincoli, vengono eseguite alcune azioni affinché il flusso termini con una o più condizioni d’uscita. \\
	\textbf{Diagramma di sequenza} & Rappresenta un diagramma che serve per descrivere uno scenario in cui le scelte o i flussi alternativi sono rappresentati tramite un pannello di "alt". In particolare servono per dare una visione specifica del flusso degli eventi a livello applicativo. \\
	\textbf{Screen mockup} & Rappresenta un prototipo dell’interfaccia che il sistema presenterà allo studente e permette al lettore di immaginare la veste grafica dell’applicazione e l’implementazione di uno specifico caso d’uso. \\
\end{tabular}
\end{table}
\clearpage
	%%%%%%%%%%%%%%%%%%%%%%%%%%%%%%%%%%%%%%%%%%%%%%%%%%%%%%%%%%%

\chapter{Analisi dei requisiti}
\label{ref:requisiti1}

%%% Il gruppo 1 scriverà qui i suoi requisiti funzionali e non funzionali %%%

\section{Requisiti funzionali}

\paragraph{Requisito 1 (sostituire con nome requisito) \\} 
Lorem ipsum dolor sit amet...

//INSERIRE QUI I REQUISITI FUNZIONALI DI TUTTI I GRUPPI

\section{Requisiti non funzionali}

Facendo riferimento al modello FURPS+ abbiamo individuato i seguenti requisiti
non funzionali.

\paragraph{Funzionalità\\} 
Il sistema sarà in grado di fruire tutte le funzionalità richieste con la massima completezza dell’applicazione.

\paragraph{Usabilità\\} 
Il sistema presenterà un’interfaccia semplice ed essenziale che garantirà la completezza e la comprensibilità per un utilizzo facile ed intuitivo per il monitoraggio della propria carriera universitaria. Al primo avvio l’utente, tramite un tutorial introduttivo, sarà guidato nell’utilizzo dell’app \emph{Studenti Unimol}.\\ 
Le informazioni presentate sullo schermo sono in grado di indirizzare l’utente verso le funzionalità a cui desidera accedere, cercando di volta in volta di isolare soltanto le informazioni necessarie.\\
L'interfaccia dell' \textit{app Studenti} mostrerà messaggi di errore con una grafica piacevole.

\paragraph{Robustezza\\} 
Ogni volta che un’operazione dell’utente sul sistema presenta un insuccesso, si invia all’utente una notifica con un messaggio di errore dando la possibilità di riprovare ad effettuare l’operazione. In caso contrario il sistema notifica il successo dell’operazione. 
Il sistema sarà anche in grado di riconoscere le parole offensive contenute all’interno di un messaggio e di non far visualizzare quest’ultimo all’interno della \emph{chat} a cui era destinato, in modo da rendere quanto più decorosa la conversazione. 

\paragraph{Performance\\} 
Il sistema dovrà ridurre al minimo i tempi di risposta in modo da rendere più piacevole e fluido l’utilizzo da parte dell’utente. Inoltre deve essere in grado di gestire più richieste da parte di diversi utenti contemporaneamente.

\paragraph{Supportabilità\\} 
Il sistema dovrà garantire semplicità nelle attività di manutenzione, risoluzione degli errori, evoluzione e aggiunta di nuove funzionalità o modifica di quelle esistenti.
Queste attività saranno garantite da aggiornamenti del \emph{software} che riguarderanno sia l’aggiunta di nuove funzionalità che l'utilizzo di nuove tecnologie per far fronte a difetti del sistema.

\paragraph{Affidabilità\\} 
Gran parte delle operazioni eseguibili sul sistema devono essere disponibili anche senza connessione \emph{Internet}. Tutto il sistema si basa sugli ultimi dati salvati nello \textit{storage}, aggiornati periodicamente.Tuttavia, il sistema richiede la connessione \emph{internet} per alcune funzionalità.

\paragraph{Sicurezza e privacy\\} 
Il sistema permetterà agli utenti del sistema l’accesso tramite credenziali. Le aree riservate e i dati sensibili saranno memorizzati sul dispositivo rispettando la normativa vigente in materia di protezione dei dati personali(\textit{GDPR}). Nella fattispecie, il sistema sarà in grado di proteggere i dati sensibili contenuti al suo interno e proteggersi dalle vulnerabilità presenti nel codice.

\paragraph{Tracciabilità\\} 
Il sistema garantirà un \emph{report} basato sull’orario di invio e conserverà informazioni delle conversazioni relative alle \emph{chat} dei corsi, al fine di poter effettuare una eventuale verifica futura.\\

\section{Pseudorequisiti e vincoli}
%%%%%%%%%%%%% 7.3.1 Implementazione %%%%%%%%%%%%%
\subsection{Implementazione}
\paragraph{}
L’app dovrà essere realizzata tramite l’utilizzo del framework \textit{Ionic}, versione 4, che comprende tecnologie come \emph{TypeScript}, \emph{Angular} e \emph{CSS}.
Per il \emph{back-end} sarà utilizzato \emph{PHP}.\\
Per il pannello di amministrazione, l'applicativo software sarà implementato con:\\
\emph{HTML}, \emph{CSS}, \emph{JavaScript} per il \emph{front-end} e \emph{PHP} per il \emph{back-end}.


%%%%%%%%%%%%% 7.3.2 Interafccia %%%%%%%%%%%%%
\subsection{Interfaccia}
\paragraph{}
L’app si appoggerà a servizi esterni offerti dall’Università utilizzando un suo sincronizzatore costruito \textit{ad hoc} per le chiamate a diversi portali.

%%%%%%%%%%%%% 7.3.3 Packaging %%%%%%%%%%%%%
\subsection{Packaging}
\paragraph{}
L’app potrà essere installata ed eseguita su tutti i dispositivi \textit{Andorid 4.4+} e \textit{iOS 9+}.

%%%%%%%%%%%%% 7.3.4 Legali %%%%%%%%%%%%%
\subsection{Legali}
\paragraph{}
L’utilizzo dell’app non comporterà il pagamento di alcuna royalty da parte dello studente.


\clearpage
	%%%%%%%%%%%%%%%%%%%%%%%%%%%%%%%%%%%%%%%%%%%%%%%%%%%%%%%%%%%

\chapter{Modello del sistema}
\label{ref:modSistema}

%%% Il gruppo 1 scriverà il suo modello del sistema. Esso dovrà includere: attori, casi d'uso (descrizione e tabella), scenari, diagrammi dei casi d'uso, diagrammi di sequenza, diagramma delle attività, screen mockups della funzionalità %%%

\section{Attori}
\label{sec:attori}
\subsection{\textit{Sync}}
\paragraph{} 
\textit{Sync} rappresenta il sincronizzatore remoto, un web service esterno che si interfaccia con una serie di servizi tra cui \textit{Esse3}, ossia il portale dello studente che offre le funzionalità da replicare nell’applicazione, \textit{Aule Unimol} ed altri servizi esterni. È un’entità esterna al sistema che si vuole realizzare, pertanto  è classificata come attore.
\begin{center}
	\includegraphics[width=0.6in]{imgs/attori/Attore-Sync.pdf}
\end{center}

\subsection{\textit{Docente}}
\paragraph{} 
\textit{Docente} è un professore dell’\textit{Università degli Studi del Molise} avente le credenziali di accesso al portale \textit{Esse3}. Il \textit{Docente} utilizzerà la nuova funzionalità per comunicare in modo rapido e semplice e anche in via ufficiosa con tutti gli studenti che frequentano un proprio corso. Egli potrà amministrare le chat secondo le sue preferenze, in particolare potrà creare o rimuovere dei canali e aggiungere o rimuovere i membri.
\begin{center}
	\includegraphics[width=0.6in]{imgs/attori/Docente.png}
\end{center}

\subsection{Studente}
\paragraph{} 
Studente regolarmente iscritto all’\textit{Università degli Studi Del Molise}, munito delle credenziali per accedere al portale dello studente \textit{Esse3}, il servizio esterno sul quale si basa l’applicazione. Lo studente utilizza l’applicazione per usufruire in maniera agevole dei servizi offerti dai portali dell’\textit{Università degli Studi del Molise} e monitorare la propria carriera universitaria. Di seguito è mostrata l'icona dello studente standard.
\begin{center}
	\includegraphics[width=0.8in]{imgs/attori/Attore-Studente.pdf}
\end{center}
Per quanto riguarda la nuova funzionalità di messaggistica, lo \textit{Studente} utilizzerà tale funzione per facilitare la comunicazione ufficiale con i docenti relativi ai corsi della Coorte di appartenenza. Potrà comunicare, inoltre, con gli altri studenti appartenenti al suo stesso anno di corso. Di seguito è mostrato lo studente della chat.
\begin{center}
	\includegraphics[width=0.8in]{imgs/attori/Studente.png}
	\label{fig:Attore: Studente chat}
\end{center}


\emph{\textbf{Amministratore:} è un Amministratore dell’Università degli Studi del Molise}, al quale sono assegnate le credenziali per poter effettuare l’accesso al pannello amministrazione.\emph{L’Amministratore} avrà il compito di gestire, mediante tale pannello, le \emph{chat} degli studenti e le \emph{chat} dei corsi dove potrà coordinare i relativi canali e gli utenti in essi presenti.
Avrà anche la possibilità di inviare notifiche personalizzate e di gestire la segnalazione dei messaggi inopportuni all’interno delle \emph{chat}.

\begin{figure}[h]
	\centering
	\includegraphics[width=0.2\textwidth]{imgs/attori/admin.png}
	\caption{Attore: Amministratore}
	\label{fig:Attore: Amministratore} 
\end{figure}



\newpage


\section{Scenari}
Inserire qui gli scenari di tutti i gruppi.

\section{Casi d'uso}
Inserire casi d'uso di tutti i gruppi.

\paragraph{Caso d'uso 1 (sostituire con nome caso d'uso) \\} 
Lorem ipsum dolor sit amet... (sostituire con descrizione caso d'uso)

\begin{table}[tb]
%\normalsize % Dimensione testo normale
\small % Dimensione testo piccola
%\footnotesize % Dimensione testo piccolissima
%\scriptsize % Dimensione del testo ulteriormente più piccola
%\caption{} % Didascalia tabella
%\label{} % Etichetta per riferimenti incrociati
\begin{tabular}{| p{\useCaseLeft} | p{\useCaseNum} | p{\useCaseTwoCol} | p{\useCaseTwoCol} |}
	\hline
	\textbf{Nome caso d'uso} & \multicolumn{3}{p{\useCaseMulticol} |}{\textbf{Login}} \\
	\hline
	\textbf{Attori partecipanti} & \multicolumn{3}{p{\useCaseMulticol} |}{Inizializzato da \textbf{Utente}.} \\
	\hline
	\textbf{Condizioni d'ingresso} & \multicolumn{3}{p{\useCaseMulticol} |}{L'utente ha cliccato sul bottone di login.} \\
	\hline
	\textbf{Flusso degli eventi} & \textbf{\#} & \textbf{Utente} & \textbf{Sistema} \\
	\hline
	\textbf{} & \textbf{1} & \textbf{} & Propone una schermata per l'inserimento dei dati necessari per il login, e-mail e password dell'utente \\
	\hline
	\textbf{} & \textbf{2} & Inserisce i dati e sottomette la richiesta & \textbf{} \\
	\hline
	\textbf{} & \textbf{3} & \textbf{} & Controlla che siano stati inseriti entrambi i campi e avvia le operazioni di visualizzazione \\
	\hline
	\textbf{Eccezioni} & \multicolumn{3}{p{\useCaseMulticol} |}{3.1 Uno o entrambi i campi sono vuoti.\newline 3.2 Le credenziali inserite non sono valide (una o entrambe).} \\
	\hline
	\textbf{Condizioni d'uscita} & \multicolumn{3}{p{\useCaseMulticol} |}{Il sistema completa la login e dà accesso all'app o, in caso contrario, visualizza un messaggio di errore se non sono stati inseriti tutti i dati obbligatori, se le credenziali non sono corrette o se si verifica un insuccesso dell'operazione.} \\
	\hline
\end{tabular}
\end{table}

\section{Diagramma dei casi d'uso}

Inserire diagramma casi d'uso di tutti i gruppi.

\begin{figure}
	\centering
	%\includegraphics[height=3in]{imgs/file-comuni-ai-gruppi/useCaseEsempio.png}
	\caption{Inserire descrizione}
	\label{fig:prova}
\end{figure}

\section{Diagramma di sequenza}

Inserire diagrammi di sequenza di tutti i gruppi.

\begin{figure}
	\centering
%	\includegraphics[height=3in,width=5in]{imgs/file-comuni-ai-gruppi/SequenceDgEsempio.png}
	\caption{Inserire descrizione}
	\label{fig:prova}
\end{figure}

\section{Diagramma delle attività}

Inserire diagrammi di activity di tutti i gruppi.

\begin{figure}
	\centering
%	\includegraphics[height=3in,width=5in]{imgs/file-comuni-ai-gruppi/ActivityDgEsempio.png}
	\caption{Inserire descrizione}
	\label{fig:prova}
\end{figure}



\clearpage

\end{document}

%%%%%%%%%%%%%%%%%%%%%%%%%%%%%%%%%%%%%%%%%%%%%%%%%%%%%%%%%%%
